\documentclass[a4paper,12pt]{article}
\usepackage[utf8]{inputenc}
\usepackage[T2A]{fontenc}
\usepackage[russian]{babel}
\usepackage{geometry}
\usepackage{setspace}
\usepackage{indentfirst}
\usepackage{graphicx}
\usepackage{booktabs}
\usepackage{longtable}
\usepackage{array}
\usepackage{multirow}
\usepackage{titlesec}

\usepackage{tabularx}
\usepackage{array}

\geometry{left=2cm, right=2cm, top=1.5cm, bottom=1.5cm}

% Отключаем гиперссылки полностью
\usepackage[hidelinks]{hyperref}

\titleformat{\section}{\normalsize\bfseries}{\thesection}{1em}{}
\titleformat{\subsection}{\normalsize\bfseries}{\thesubsection}{1em}{}
\titleformat{\subsubsection}{\normalsize\bfseries}{\thesubsubsection}{1em}{}

\renewcommand{\contentsname}{Содержание}

\begin{document}

% Титульный лист
\begin{titlepage}
\noindent
\begin{tabular*}{\textwidth}{@{}p{0.48\textwidth}@{\extracolsep{\fill}}p{0.48\textwidth}@{}}
	\centering\textbf{УТВЕРЖДАЮ} & \centering\textbf{УТВЕРЖДАЮ} \tabularnewline
	& \tabularnewline
	& \tabularnewline
	& \tabularnewline
	\rule{5cm}{0.4pt} /Кравченко О.Ю./ & \rule{5cm}{0.4pt} /Кравченко О.Ю./ \tabularnewline
	\centering подпись & \centering подпись \tabularnewline
	& \tabularnewline
	& \tabularnewline
	& \tabularnewline
	«\rule{1.5cm}{0.4pt}» \rule{4cm}{0.4pt} 2026 год & «\rule{1.5cm}{0.4pt}» \rule{4cm}{0.4pt} 2026 год \tabularnewline
	\centering дата & \centering дата \tabularnewline
\end{tabular*}

\vspace{0.5cm}

\begin{center}
	\textbf{ТЕХНИЧЕСКОЕ ЗАДАНИЕ} \\
	\textbf{НА СОЗДАНИЕ САЙТА ДЛЯ СРАВНЕНИЯ ЦЕН} \\
	\textbf{НА ТОВАРЫ МЕЖДУ МАРКЕТПЛЕЙСАМИ} \\
	\textbf{OZON И WILDBERRIES} \\
\end{center}

\vspace{0.5cm}

% Таблица с информацией - делаем компактнее
\footnotesize
\noindent
\begin{tabular}{|p{0.32\textwidth}|p{0.28\textwidth}|p{0.18\textwidth}|p{0.18\textwidth}|}
	\hline
	Заказчик, утверждающий ТЗ & Кравченко О.Ю. +7 (908) 517 50 43 & Аналитик-системный архитектор & ----- \\
	\hline
	Заказчик, согласующий ТЗ & Кравченко О.Ю. +7 (908) 517 50 43 & Дата начала подготовки & 12:00 26.10.2025 \\
	\hline
	Шифр проекта & «Price Hunter Pro» & Дата текущей версии & 23:54 29.03.2026 \\
	\hline
	Сотрудники Заказчика, принимавшие участие в проработке технического задания & ----- & Версия документа & 1 \\
	\hline
	Готовность & 90\% & Автор изменений &-----\\
	\hline
\end{tabular}

\vspace{0.5cm}

\vfill

\noindent
\hspace{2.5cm} \textbf{СОГЛАСОВАНО}

\vspace{1.5cm}

\noindent
\begin{tabular}{@{}p{0.6\textwidth}@{}}
	\rule{6cm}{0.4pt} /Кравченко О.Ю./ \\
	\hspace{2.4cm} \small подпись \\[2cm]
	«\rule{1.5cm}{0.4pt}» \rule{5cm}{0.4pt} 2026 год \\
	\centering{\small дата} \\
\end{tabular}

    \vfill
    \begin{center}
        Ростов-на-Дону, 2026
    \end{center}
\end{titlepage}

% Содержание
\newpage
\tableofcontents
\newpage


\textbf{ТЕХНИЧЕСКОЕ ЗАДАНИЕ САЙТА ДЛЯ СРАВНЕНИЯ ЦЕН НА ТОВАРЫ МЕЖДУ МАРКЕТПЛЕЙСАМИ OZON И WILDBERRIES}


\section{ГЛОССАРИЙ}

\begin{description}
	\item[БД] — База данных — совокупность структурированных данных, хранящихся в компьютере.
	
	\item[Бекэнд] — Программный комплекс, обрабатывающий информацию на стороне сервера; внутренняя реализация, не показываемая пользователю.
	
	\item[Веб-приложение] — Клиент-серверное приложение, в котором клиент взаимодействует с веб-сервером при помощи браузера.
	
	\item[Интернет] — Глобальная информационно-коммуникационная сеть.
	
	\item[Маркетплейс] — Торговая площадка в интернете, на которой продавцы предлагают свои товары, а покупатели могут их приобрести (в данном проекте — Ozon и Wildberries).
	
	\item[Парсинг] — Автоматический сбор и структурирование информации с веб-страниц с помощью специализированных программ.
	
	\item[ПО] — Программное обеспечение.
	
	\item[Сайт] — Программный комплекс, обрабатывающий информацию на стороне клиента и обеспечивающий доступ к ней через интернет.
	
	\item[ТЗ] — Техническое задание — документ, определяющий требования к создаваемой системе.
	
	\item[Фронтэнд] — Презентационная часть веб-приложений, информационной или программной системы, её пользовательский интерфейс и связанные с ним компоненты.
	
	\item[API] — Интерфейс программирования приложений, позволяющий одной программе взаимодействовать с другой для обмена данными и функциями.
	
	\item[Кэширование] — Промежуточное хранение данных в быстродействующем хранилище для ускорения доступа к ним при повторных запросах.
	
	\item[Партнерская ссылка] — Специальная ссылка, содержащая идентификатор партнера, позволяющая отслеживать переходы пользователей на сайт маркетплейса и получать комиссию за совершенные покупки.
\end{description}

\newpage

\section{ОБЩИЕ СВЕДЕНИЯ}

\subsection{Полное наименование системы и ее условное обозначение}
Полным наименованием является «Price Hunter Pro: сервис сравнения цен
между маркетплейсами Ozon и Wildberries».

\subsection{Шифр проекта}
Шифрованным обозначением проекта является «Price Hunter Pro».

\subsection{Заказчик, исполнитель и вовлечённые в проект лица}
В создание Price Hunter Pro вовлечены физические лица, перечисленные в таблице 2.

\begin{table}[h]
    \centering
    \caption{Вовлечённые в создание Price Hunter Pro физические лица}
    \begin{tabular}{|p{4cm}|p{5cm}|p{4cm}|}
        \hline
        \textbf{Фамилия, имя, отчество} & \textbf{Роль в проекте} & \textbf{Контакт} \\
        \hline
        Корепина Анна Алексеевна & Веб-дизайнер & +79612832135 \\
        \hline
        Коростиенко Максим Александрович & Backend разработчик & +79718302754 \\
        \hline
        Панировский Иван Владимирович & Frontend разработчик & +79381346115 \\
        \hline
    \end{tabular}
\end{table}

\subsection{Перечень документов и консультаций, на основании которых создаётся Price Hunter Pro}

\subsubsection{Перечень документов, на основании которых создаётся Price Hunter Pro}
\begin{enumerate}
    \item Федеральный закон от 27.07.2006 N 149-ФЗ «Об информации, информационных технологиях и о защите информации».
    \item Федеральный закон от 27.07.2006 N 152-ФЗ «О персональных данных».
    \item Гражданский кодекс Российской Федерации (часть четвертая) от 18.12.2006 N 230-ФЗ — в части соблюдения авторских прав и интеллектуальной собственности.
    \item Правила использования сайтов маркетплейсов Ozon и Wildberries (условия партнерских соглашений и требования к автоматизированному сбору данных).
\end{enumerate}

\subsubsection{Перечень консультаций, на основании которых создаётся Price Hunter Pro}
\begin{enumerate}
    \item Консультация с потенциальными пользователями сервиса в период с 01.10.2025 по 15.10.2025. В ходе опроса выявлена потребность в быстром сравнении цен на идентичные товары между двумя крупнейшими маркетплейсами.
    \item Анализ существующих решений на рынке (агрегаторы цен, аналогичные сервисы). Выявлены их сильные и слабые стороны, сформированы требования к уникальному функционалу.
    \item Изучение документации партнерских программ Ozon и Wildberries, проведенное 20.10.2025. Определены технические требования к формированию партнерских ссылок и ограничения по частоте запросов к данным.
    %\item Консультация со специалистом по парсингу данных 25.10.2025. Получены рекомендации по выбору инструментов для сбора данных и обходу блокировок при соблюдении правил robots.txt.
    \item Технический анализ структуры сайтов Ozon и Wildberries, проведенный 01.11.2025. Выявлены ключевые параметры товаров для сопоставления (название, артикул, характеристики).
\end{enumerate}

\subsection{Плановые сроки начала и окончания работы по созданию сайта}
Плановый срок начала работ создания Price Hunter Pro — 01.10.2025. 

Плановый срок окончания работ создания Price Hunter Pro — 31.05.2026.

\subsection{Сведения об источниках и порядке финансирования работ}
Финансирование работ осуществляется за счёт собственных средств разработчиков. Начальные инвестиции отсутствуют, используются технологии open-source. Эксплуатационные расходы предполагают оплату хостинга и доменного имени.

\subsection{Порядок оформления и предъявления заказчику результатов работ по созданию сайта (его частей), а также оплаты заказчиком работы исполнителя}
ТЗ оформляется в электронном виде в формате, совместимом с LaTeX и PDF. По завершению разработки заказчику предоставляются:
\begin{itemize}
    \item Демонстрация работоспособности системы;
    \item Полный пакет документации (ТЗ, программа и методики испытаний, пояснительная записка, руководство пользователя);
    \item Исходный код проекта с комментариями.
\end{itemize}
Приемка работ осуществляется после успешного прохождения всех видов испытаний и подписания акта приемки-передачи. Оплата работ исполнителю не предусмотрена в связи с инициативным характером разработки в учебных целях.

\newpage
\section{ЦЕЛИ И НАЗНАЧЕНИЕ СОЗДАНИЯ САЙТА СРАВНЕНИЯ ЦЕН}

\subsection{Цели создания сайта}
Целью создания Price Hunter Pro является автоматизация процесса сравнения цен на товары между маркетплейсами Ozon и Wildberries. Основные цели:
\begin{itemize}
    \item Сокращение времени поиска оптимальных цен для пользователей;
    \item Предоставление удобного инструмента для анализа рыночных цен;
    \item Создание платформы для монетизации через партнерские программы;
    \item Повышение эффективности онлайн-покупок.
\end{itemize}

\subsection{Наименования и требуемые значения технических показателей сайта}
Price Hunter Pro должен обеспечивать следующие технические показатели:
\begin{itemize}
    \item Время ответа системы: не более 5 секунд при стандартном поисковом запросе.
    \item Поддержка одновременной работы: не менее 100 одновременных пользователей.
    \item Доступность системы: не менее 99.5\% времени (за исключением плановых технических работ).
    \item Объем базы данных: возможность хранения информации о 10 000+ уникальных товарах.
\end{itemize}

\subsection{Наименования и требуемые значения технологических показателей сайта}
Price Hunter Pro должен обеспечивать следующие технологические показатели:
\begin{itemize}
    \item Прием и обработка поисковых запросов от пользователей через веб-интерфейс;
    \item Параллельный сбор данных с двух маркетплейсов с использованием методов парсинга;
    \item Автоматическое сопоставление идентичных товаров по ключевым параметрам (название, артикул, бренд, основные характеристики);
    \item Кэширование результатов популярных запросов для снижения нагрузки на маркетплейсы и ускорения работы;
    \item Формирование пользовательских карточек товаров с отображением цен и ссылок на оба маркетплейса;
    \item Обработка ошибок и исключительных ситуаций (отсутствие товара, недоступность сайта, блокировка парсера).
\end{itemize}

\subsection{Наименования и требуемые значения \\производственно-экономических показателей сайта}
Price Hunter Pro должен обеспечивать следующие производственно-экономические показатели:
\begin{itemize}
    \item Минимизация затрат на разработку за счёт использования open-source технологий;
    \item Потенциальная окупаемость за счёт монетизации через партнерские программы Ozon и Wildberries;
    \item Снижение трудозатрат пользователей на поиск товаров по оптимальным ценам (экономия времени).
\end{itemize}

\subsection{Наименования и требуемые значения других показателей сайта}

\subsubsection{Техническое обеспечение}

\textbf{Серверная часть:}
\begin{itemize}
    \item Язык программирования: Python 3.9+;
    \item Веб-фреймворк: FastAPI;
    \item Парсинг данных: Playwright (библиотека Camoufox).
\end{itemize}

\textbf{Клиентская часть:}
\begin{itemize}
    \item HTML5, CSS3, JavaScript (ES6+);
    \item Метод запросов: Fetch API;
    \item UI-библиотека: Bootstrap 5.x;
    \item Хранение данных: LocalStorage API.
\end{itemize}

\textbf{Аппаратные требования:}
\begin{itemize}
    \item Интернет-соединение.
\end{itemize}

\subsubsection{Ограничения и запрещенные функции}
В целях соблюдения законодательства и правил маркетплейсов, система НЕ должна:
\begin{itemize}
    \item Осуществлять прямое совершение покупок через приложение;
    \item Хранить персональные данные пользователей (логины, пароли, платежную информацию, историю покупок);
    \item Предоставлять возможность оставлять отзывы о товарах или продавцах;
    \item Осуществлять модификацию цен или условий продажи товаров;
    \item Предоставлять услуги доставки товаров или гарантии на них;
    \item Создавать пользовательские аккаунты и личные кабинеты с длительным хранением истории;
    \item Осуществлять автоматическое обновление цен чаще чем 1 раз в 2 часа;
    \item Собирать данные, не относящиеся к сравнению цен (контактные данные продавцов, персональные отзывы пользователей, историю заказов);
    \item Нести ответственность за качество товаров;
    \item Использовать полученные данные для создания конкурентной базы или перепродажи третьим лицам.
\end{itemize}

\subsection{Назначение сайта}
Price Hunter Pro предназначен для автоматизации процесса сравнения цен на товары между маркетплейсами Ozon и Wildberries, предоставляя пользователям актуальную информацию о стоимости товаров на обеих платформах в одном интерфейсе.

В рамках работы Price Hunter Pro достигаются следующие улучшения:
\begin{itemize}
    \item Сокращается время пользователя на поиск выгодных предложений (исключается необходимость вручную открывать оба сайта и искать один и тот же товар);
    \item Предоставляется инструмент для быстрого анализа цен на разных платформах с наглядным отображением разницы в стоимости;
    \item Создается основа для монетизации через партнерские программы (комиссия за переходы и покупки);
    \item Повышается эффективность онлайн-покупок за счёт принятия обоснованных ценовых решений;
    \item Обеспечивается прозрачность ценообразования на рынке электронной коммерции;
    \item Стимулируется здоровая конкуренция между маркетплейсами через информирование покупателей.
\end{itemize}
\newpage
\section{ХАРАКТЕРИСТИКА ОБЪЕКТОВ}

\subsection{Основные сведения об объекте автоматизации или ссылки на документы, содержащие такие сведения}

Подготовленный сайт должен выполнять следующие функции:

\begin{enumerate}
	\item Взаимодействовать с маркетплейсами Ozon и Wildberries посредством парсинга данных в соответствии с партнерскими соглашениями.
	
	\item Иметь основной тип пользователей - неавторизованные пользователи (гостевой доступ)

	
	\item Иметь возможность поиска и сравнения цен на идентичные товары по определённым параметрам:
	\begin{itemize}
		\item Наименование товара
		\item Бренд
		\item Основные характеристики
	\end{itemize}
	
	\item Должен предоставлять возможность формирования партнёрских ссылок для перехода на маркетплейсы.
\end{enumerate}

\subsection{Сведения об условиях эксплуатации сайта и характеристиках окружающей среды}

Сайт должен быть спроектирован таким образом, чтобы одновременно в любое время суток им могли пользоваться \textbf{100 (сто) человек} одновременно с возможностью дополнительного расширения до \textbf{1000 (одной тысячи) человек} одновременно.

Доступ на сайт осуществляется через интернет с любых устройств (ПК, планшеты, мобильные телефоны) при наличии современного веб-браузера.

Все данные (кроме кэшированных результатов поиска) обязательно должны храниться в базе данных на сервере хостинг-провайдера. Персональные данные пользователей не хранятся в соответствии с требованиями безопасности.

Время ответа системы не должно превышать \textbf{5 секунд} при стандартном поисковом запросе.

Доступность системы должна составлять не менее \textbf{99.5\%} времени (за исключением плановых технических работ).

\newpage
\section{ТРЕБОВАНИЯ К СИСТЕМЕ}

\subsection{Требования к системе в целом}

\subsubsection{Требования к структуре и функционированию системы}

\textbf{Перечень подсистем, их назначение и основные характеристики, требования к числу уровней иерархии и степени централизации системы}

Система должна включать в себя следующие подсистемы:

\textbf{I. Пользовательский интерфейс (Frontend):}
\begin{enumerate}
	\item Страница поиска товаров
	\item Страница результатов сравнения цен
	\item Страница детальной информации о товаре
	\item Адаптивная вёрстка для мобильных устройств
\end{enumerate}

\textbf{II. Серверная часть (Backend):}
\begin{enumerate}
	\item Модуль приёма и обработки поисковых запросов
	\item Модуль парсинга данных с Ozon
	\item Модуль парсинга данных с Wildberries
	\item Модуль сопоставления товаров
	\item Модуль кэширования результатов
\end{enumerate}

\textbf{III. База данных:}
\begin{enumerate}
	\item Хранение информации о товарах (наименование, артикул, характеристики)
	\item Хранение кэшированных данных о ценах
	\item Хранение статистики поисковых запросов
\end{enumerate}

\textbf{IV. Система обработки ошибок:}
\begin{enumerate}
	\item Обработка недоступности маркетплейсов
	\item Обработка отсутствия товаров
	\item Обработка блокировок парсера
	\item Логирование ошибок
\end{enumerate}

\textbf{V. Безопасность:}
\begin{enumerate}
	\item Защита от DDoS-атак
	\item Ограничение частоты запросов
	\item Валидация входных данных
	\item Соответствие ФЗ-152 «О персональных данных»
\end{enumerate}

\subsection{Требования к способам и средствам связи для информационного обмена между компонентами системы}

Весь обмен информацией происходит через Интернет.

\textbf{Взаимодействие Frontend-Backend:}
\begin{itemize}
	\item Связь осуществляется посредством отправки HTTP/HTTPS запросов (GET, POST)
	\item Формат обмена данными: JSON
	\item API сайта должно соответствовать протоколу REST
\end{itemize}

\textbf{Взаимодействие Backend-Маркетплейсы:}
\begin{itemize}
	\item Получение информации от Ozon происходит посредством парсинга с использованием библиотеки Playwright (Camoufox)
	\item Получение информации от Wildberries происходит посредством парсинга с использованием библиотеки Playwright (Camoufox)
	\item Частота обновления данных: не чаще 1 раза в 2 часа
\end{itemize}

\textbf{Взаимодействие Backend-База данных:}
\begin{itemize}
	\item Подключение к СУБД осуществляется через ORM
	\item Поддерживается работа с PostgreSQL/MySQL
\end{itemize}

\textbf{Требования к характеристикам взаимосвязей создаваемой системы со смежными системами, требования к её совместимости, в том числе указания о способах обмена информацией:}

Сайт должен быть совместим с:
\begin{itemize}
	\item Браузерами: Chrome 90+, Firefox 88+, Safari 14+, Edge 90+
	\item Мобильными браузерами: Chrome Mobile, Safari iOS
\end{itemize}

Обмен информацией с маркетплейсами происходит автоматически через парсинг в соответствии с их правилами использования.

\subsection{Требования к режимам функционирования системы}

Сайт Price Hunter Pro должен функционировать ежедневно в режиме круглосуточного доступа (24 часа в сутки, 7 дней в неделю).

\subsection{Требования по диагностированию системы}

Сайт считается работоспособным, если одновременно весь функционал доступен \textbf{100 (ста) пользователям} одновременно, в зависимости от ролей. Сайт должен быть доступен на русском языке.

Система должна обеспечивать:
\begin{itemize}
	\item Автоматический мониторинг доступности маркетплейсов Ozon и Wildberries;
	\item Контроль времени ответа системы (не более 5 секунд);
	\item Отслеживание успешности парсинга данных;
	\item Логирование всех критических ошибок и сбоев;
	\item Автоматическое оповещение разработчиков о критических сбоях системы.
\end{itemize}

\subsection{Перспективы развития, модернизации системы}

При необходимости сайт Price Hunter Pro может быть масштабирован до \textbf{1000 (одной тысячи) пользователей}, которым доступен одновременно весь функционал системы, в зависимости от ролей.

Должна быть возможность:
\begin{itemize}
	\item Добавления поддержки дополнительных маркетплейсов (Яндекс.Маркет, AliExpress Russia, СберМегаМаркет);
	\item Перевода системы на другие языки (английский, казахский, белорусский);
	\item Расширения функционала за счёт внедрения системы уведомлений об изменении цен;
	\item Создания мобильного приложения на базе существующего API;
	\item Внедрения системы аналитики и статистики для пользователей;
	\item Добавления функции отслеживания истории изменения цен на товары;
	\item Интеграции с социальными сетями для возможности делиться находками;
	\item Создания API для сторонних разработчиков.
\end{itemize}

Также коробочную версию можно разворачивать:
\begin{itemize}
	\item В коммерческих организациях для внутреннего использования (корпоративные закупки);
	\item В государственных учреждениях для мониторинга цен на закупаемые товары.
\end{itemize}
\end{document}